%% ma: L12-B03-0005
%% lop: 12
%% chuong: 1
%% bai: 3
%% dang: 12.3.4
%% loai: TN4
%% muc: TH
%% dap_an: "D"
%% nguon: "(NBV)-ĐỀ SỐ 6-MỖI NGÀY 1 ĐỀ THI 2025.pdf (D06, MỖI NGÀY 1 ĐỀ - Đề số 6), Phần I Câu 5"
%% kien_thuc: "Bài 3 (tiệm cận đứng, tiệm cận ngang của đồ thị hàm phân thức)"
%% trang_thai: cho-duyet
%% ly_do: "Dạng toán mới đề xuất 12.3.4 (đếm số đường tiệm cận của đồ thị hàm phân thức bậc hai trên bậc hai), chờ giáo viên duyệt"
%% ghi_chu: "Hình dựng lại bằng plot từ hàm y=-x^2/(x^2-1) (tiệm cận đứng x=-1, x=1; tiệm cận ngang y=-1). Đáp án gốc D đúng"
\begin{ex}
Cho hàm số $y=f(x)=\dfrac{-x^2}{(x-1)(x+1)}$ có đồ thị như hình vẽ.
\begin{center}
\begin{tikzpicture}[x=1cm,y=.7cm]
  \begin{scope}
    \clip (-3.4,-4.2) rectangle (3.4,4.4);
    \draw[thick,domain=-.92:.92,samples=80,smooth] plot (\x,{(\x*\x)/(1-\x*\x)});
    \draw[thick,domain=1.1:3.3,samples=80,smooth] plot (\x,{-(\x*\x)/(\x*\x-1)});
    \draw[thick,domain=-3.3:-1.1,samples=80,smooth] plot (\x,{-(\x*\x)/(\x*\x-1)});
  \end{scope}
  \draw[phu] (-1,-4.2)--(-1,4.4);
  \draw[phu] (1,-4.2)--(1,4.4);
  \draw[phu] (-3.4,-1)--(3.4,-1);
  \draw[truc] (-3.5,0)--(3.6,0) node[below]{$x$};
  \draw[truc] (0,-4.3)--(0,4.6) node[left]{$y$};
  \foreach \p in {-1,1} \draw (\p,.08)--(\p,-.08) node[below,font=\footnotesize]{$\p$};
  \draw (.06,-1)--(-.06,-1) node[left,font=\footnotesize]{$-1$};
  \path (0,0) node[below left,font=\footnotesize]{$O$};
\end{tikzpicture}
\end{center}
Số đường tiệm cận của đồ thị hàm số đã cho là
\choice{$0$}{$1$}{$2$}{\True $3$}
\loigiai{Ta có $\lim\limits_{x\to\pm\infty}\dfrac{-x^2}{x^2-1}=-1$ nên đồ thị có tiệm cận ngang $y=-1$. Vì $\lim\limits_{x\to1^+}f(x)=-\infty$ và $\lim\limits_{x\to-1^+}f(x)=+\infty$ nên đồ thị có hai tiệm cận đứng $x=1$ và $x=-1$. Vậy đồ thị có $3$ đường tiệm cận.}
\end{ex}
